La transformation numérique des organisations a profondément reconfiguré les pratiques
de gestion de projets. Les équipes de développement logiciel évoluent désormais dans
des environnements où la complexité, la rapidité et la collaboration distribuée exigent
des outils à la hauteur de leurs ambitions. Les plateformes existantes — Jira, Trello,
Asana — répondent partiellement à ces besoins, mais aucune n'intègre nativement une
couche d'intelligence artificielle capable d'assister les décisions opérationnelles.

\bigskip
C'est dans ce contexte que s'inscrit ce projet de fin d'études, réalisé au sein de
l'entreprise \textbf{TAC-TIC} (Groupe New Ways, Technoparc El Ghazela, Ariana,
Tunisie). La mission confiée : concevoir et développer \textbf{TAC-TIC IntelliPM},
une plateforme de gestion de projets intelligente combinant architecture distribuée,
sécurité multicouche et moteur IA hybride.

\bigskip
\noindent\textbf{Organisation du rapport}

\medskip
Ce document est structuré en \textbf{six chapitres} :

\begin{description}[leftmargin=3.5cm, style=nextline]
  \item[Chapitre I] \textit{Cadre général du projet} — présentation de TAC-TIC,
        contexte, problématique, étude de l'existant, méthodologie Scrum et
        planification des six sprints.

  \item[Chapitre II] \textit{Planification et Architecture} — analyse des besoins
        (acteurs, rôles, RBAC), spécification fonctionnelle et non fonctionnelle,
        modélisation UML, product backlog (34 user stories), conception architecturale
        distribuée et modèle de données.

  \item[Chapitre III] \textit{Réalisation des modules fondamentaux} — trois premiers
        sprints : authentification Sanctum et RBAC Spatie (Sprint~1), gestion des
        projets, équipes et invitations (Sprint~2), gestion des tâches avec Kanban
        drag-and-drop, dépendances BFS et Gantt CPM (Sprint~3).

  \item[Chapitre IV] \textit{Réalisation des modules avancés} — Sprint~4 (notifications
        temps réel WebSocket, dashboards multi-rôles, calendrier, recherche globale)
        et Sprint~5 (moteur IA hybride : recommandation ML via FastAPI, estimation
        de durée par heuristiques, détection de risques, assistant LLM).

  \item[Chapitre V] \textit{Administration et pilotage} — Sprint~6 : système RBAC
        entièrement personnalisable (rôles, permissions atomiques, matrice), journal
        d'audit enterprise (valeurs avant/après, sévérité, IP, export PDF, broadcast
        WebSocket) et tableau de bord analytique (KPIs, détection d'anomalies).

  \item[Chapitre VI] \textit{Tests, Validation et Déploiement} — stratégie de tests
        fonctionnels, analyse SonarCloud (Quality Gate passé), conteneurisation Docker,
        pipeline CI/CD GitHub Actions et déploiement Azure.
\end{description}

\medskip
Le rapport se clôt par une \textbf{conclusion générale} récapitulant les résultats,
les limites identifiées et les perspectives d'évolution.
